\section{The Question Beneath Every Other Question}

Can a creature say without reserve, ``Thy will be done,'' and still act, choose, oppose injustice, repent of sin, and answer for what he has done? If God governs all things, is human freedom only an appearance? If a human being is genuinely free, must some corner of reality stand outside Providence? If suffering is received from the Father's hand, does that excuse the hand of an abuser? If charity says, ``do what you will,'' has law disappeared---or has the will itself been made new?

These are not separate puzzles. They are one question viewed from above and below: \emph{How does the causality of the Creator relate to the causality of creatures?} The spiritual life goes astray when it imagines God and the human person as two agents competing on one level. The metaphysical mistake then becomes a moral and pastoral wound. One person protects freedom by making Providence distant; another protects Providence by making the person passive. One makes surrender an anxious effort to decode a hidden script; another makes it acquiescence before sin. Catholic faith refuses all four reductions.

\articlethesis{God's rule and creaturely freedom are not rival powers. Catholic doctrine requires that his one eternal will infallibly govern all things without causing moral evil or destroying created freedom. The Thomistic synthesis employed here explains that God gives the positive reality and action by which a creature acts, while the sinful defect arises from the creature's failure to order that act to the good. No sin escapes or defeats God's wisdom, justice, or mercy. Trustful surrender therefore means doing the good entrusted to us while placing every uncontrollable result under the Father Almighty revealed in Christ.}

\subsection{Reader, scope, sources, and limits}

This is a systematic and spiritual-theological article for serious lay study. It seeks the doctrinal architecture beneath a practice of surrender: Scripture supplies the grammar; councils and the \emph{Catechism of the Catholic Church} mark governing boundaries; Fathers show the mystery prayed and contested; Aquinas gives a causal and moral analysis; Jean-Baptiste Saint-Jure and St.~Claude de la Colombière supply a compact ascetical school. A short clarification of Augustine's later resemblance to Aleister Crowley's formulas states only the history the sources establish and keeps theological comparison distinct from textual genealogy.

The article does not decide the permitted disputes between Thomists, Molinists, and other Catholic schools over how divine foreknowledge, providential concurrence, grace, and free consent are coordinated. Nor is it a diagnosis, a direction for one endangered person, or permission to remain within violence. Flight, defense, reporting, treatment, lament, prosecution, and the protection of dependents can be works of justice and charity. A sufferer who cannot yet feel peace has not thereby failed God. The Cross permits lament; surrender concerns the will's filial adherence amid feeling, not the production of a serene mood.

\begin{threestable}{Source class}{Function in this article}{Governing use}
Sacred Scripture & Inspired, normative witness read within the Church; the article's exegesis remains identified as exegesis. & Joseph, Job, Gethsemane, the Passion, Pauline freedom, and Johannine charity.\\
Magisterium & Defined doctrine and other authoritative teaching, calibrated by document and locus rather than flattened into one rank. & Omnipotence, universal Providence, God's innocence of moral evil, authentic freedom, moral acts, and the priority of grace.\\
Patristic witness & Evidence of received Christian grammar and exegesis; breadth is not a mechanical authority count. & Irenaeus, Gregory of Nyssa, Chrysostom, Augustine, Gregory the Great, and John Damascene.\\
Spiritual witness & Saintly or ascetical teaching received for prayer and practice, not treated as a technical definition or magisterial act. & Saint-Jure and La Colombière on duties, trial, petition, and abandonment.\\
Thomistic theology & Aquinas's specifically theological causal articulation, employed within the doctrinal boundaries above and not treated as definition where permitted Catholic schools disagree. & Primary and secondary causality, willed permission, voluntary action, the act and its defect, and essentially ordered causes.\\
Editorial synthesis & This project's source-grounded integration or prudential application, attributable to the article alone. & The paired-regress diagram, the agency / culpability map, and the bounded theological comparison of Augustine and Crowley.\\
\end{threestable}

The full audit identifies exact loci, online witnesses, rejected claims, and outstanding review in the repository source record. This publication is therefore a \emph{source-audited working article}; it has not received independent theological review or ecclesiastical approval.

\subsection{The short treatise received---and bounded}

The English volume \emph{Trustful Surrender to Divine Providence} is a devotional anthology. Its first part extracts Jean-Baptiste Saint-Jure's \emph{The Knowledge and Love of Our Lord Jesus Christ}; its second presents an extract from St.~Claude de la Colombière. Its governing spiritual claim begins with a decisive fence: ``apart from sin,'' no event escapes God's will or permission. Properly understood, that is not a division between strong action and weak tolerance: what God does not produce as good can occur only under a knowing permission enclosed within his omnipotent government. The anthology then teaches conformity through ordinary vexation, illness, loss, family duty, interior trial, prayer, and hope.

This article receives the anthology as a teacher of spiritual vision, not as a technical manual whose every compressed phrase is a doctrinal formula. Its opening exception for sin is decisive: the malice belongs to the wrongdoer, never to God. Thomistic theology can say that a physical loss is willed incidentally when it is inseparable from a good God positively wills; the freely offered Passion is the unique saving summit of this truth. But another person's sin never authorizes us to tell a victim that God positively desired the injury or to claim knowledge of his hidden reason. What can always be confessed is that neither the injury nor its consequences escaped his government. The anthology's best practical order is equally important: first keep God's commandments, obey legitimate authority, and perform the duties of one's state; then receive what Providence decides beyond one's power. Its section on recourse to prayer makes surrender active and petitionary rather than mute.

Yet devotional intensity needs dogmatic grammar. When the anthology calls God the cause of an illness or says that God is ``responsible'' for an injury, the language must be read through the distinction it itself announces and through the Church's explicit teaching: ``God is in no way, directly or indirectly, the cause of moral evil'' (CCC 311). God never supplies the malice; no created event falls outside his government. That qualification will govern every page that follows.

\begin{massbox}{The foundational rule}
Receive the \emph{event} as held within Providence; judge the \emph{act} by truth and moral law; oppose the \emph{evil} when charity assigns that duty; forgive the \emph{person} without calling the sin good; entrust the \emph{outcome} to God. None of these verbs cancels another.
\end{massbox}
